\documentclass[a4paper]{scrartcl}

\newcommand{\titlename}{Vergleich und Verbesserung von heuristischer und graphbasierter Algorithmen für das Spiel Snake}
\newcommand{\authorname}{Moritz Köhler}
\newcommand{\boardsize}{30\,$\times$\,30}

\usepackage[
  pdftitle={Disposition zur Studienarbeit: \titlename},
  pdfauthor={\authorname},
  pdfsubject={Disposition zur Studienarbeit},
  pdfkeywords={Snake, Algorithmen, A*, Hamiltonian, Greedy, Replay-Format, State-Saver, Game AI},
  colorlinks=true,
  linkcolor=black,
  urlcolor=cyan,
  citecolor=black
]{hyperref}
\usepackage{fontenc}
\usepackage{amssymb}

\title{\Large Exposé zur Studienarbeit\\\LARGE \titlename}
\author{\authorname}
\date{\today}

\begin{document}

\maketitle

\section{Kurzbeschreibung der Arbeit}
% Eigentlich n x n statt 30 x 30
Ziel der Studienarbeit ist die Entwicklung, der Vergleich und die iterative Verbesserung mehrerer Algorithmen für das Spiel "Snake". Im Mittelpunkt steht ein \boardsize-Feld, auf dem verschiedene Ansätze wie A*, Hamiltonian-basierte Strategien, Greedy-Heuristiken und hybride Varianten getestet werden. Die Arbeit soll zeigen, welche Verfahren unter identischen Bedingungen die beste Spielstärke, Stabilität und Laufzeit erreichen.

Die Spielstärke hat zwei Metriken. Die Win Rate (also wie oft das Spiel erfolgreich abgeschlossen wird) und die Anzahl an Schritten bis zum Spielende. Stabilität wird durch die Varianz der Ergebnisse über viele Partien gemessen, während die Laufzeit die Effizienz der Algorithmen bewertet. Die Arbeit soll nicht nur die besten Ansätze identifizieren, sondern auch deren Schwächen analysieren und gezielt Verbesserungen ableiten.

Für die Entwicklung ist ein State Management und Replay System vorgesehen, das Spielzustände, Züge und Metadaten speichert und später reproduzierbar wieder abspielt. Dadurch lassen sich einzelne Partien analysieren und strategische Fehler sichtbar machen.

\section{Gliederung und Zeitplan}
\begin{enumerate}
\item Einordnung des Problems und kurze Übersicht vorhandener Snake-Ansätze
\item Entwurf des State Managements und Replay-Systems
\item Implementierung erster Basisalgorithmen für ein \boardsize-Feld
\item Test und Vergleich der Algorithmen anhand von Simulationen und Metriken
\item Analyse der Schwächen und Ableitung verbesserter Strategien
\item Implementierung der verbesserten Strategien und erneuter Vergleich
\item Dokumentation der Ergebnisse und Bewertung der Ansätze
\item Optionale Erweiterungen:
    \begin{itemize}
        \item Erweiterung der Feldgröße auf n\,$\times$\,m ($n, m, \in \mathbb{N}$) und Analyse der Skalierbarkeit mittels O-Notation
        \item Vergleich mit weiteren Algorithmen oder KI-Ansätzen
        \item Vergleich zu Menschen gespielten Snake Spielen ohne die Notwendige Reaktionszeit für Vergleichbarkeit.
    \end{itemize}
\end{enumerate}

\section{Grundlegende Literatur}
% "Quellen Arbeit ist sehr sauber" - Felix
Für die algorithmische Grundlage sind A* sowie die Heuristiksuche zentral~\cite{Hart1968,Edelkamp2011}. Für Hamiltonische Ansätze sind Arbeiten zu Hamilton-Pfaden in Gittergraphen relevant~\cite{Itai1982,Chvatal1985}. Für die Entwicklung von KI basierten Ansätzen sind weitere Quellen hilfreich~\cite{Justesen2019,Togelius2009}. Als aktuelle Referenzimplementierungen und Inspirationsquellen dienen zudem vorhandenes Snake-Projekt~\cite{alphaPhoenix}.

\begin{thebibliography}{7}
% https://www.cs.auckland.ac.nz/courses/compsci709s2c/resources/Mike.d/astarNilsson.pdf
% https://ieeexplore.ieee.org/document/4082128
% https://doi.org/10.1109/TSSC.1968.300136
\bibitem{Hart1968} P.~E. Hart, N.~J. Nilsson, and B. Raphael. A formal basis for the heuristic determination of minimum cost paths. \textit{IEEE Transactions on Systems Science and Cybernetics}, 4(2):100--107, 1968 \href{https://doi.org/10.1109/TSSC.1968.300136}{https://doi.org/10.1109/TSSC.1968.300136}.
% https://doi.org/10.1016/C2009-0-16511-X
\bibitem{Edelkamp2011} S. Edelkamp and S. Schroedl. \textit{Heuristic Search: Theory and Applications}. Morgan Kaufmann, 2011 \href{https://doi.org/10.1016/C2009-0-16511-X}{https://doi.org/10.1016/C2009-0-16511-X}.
% https://www.ibr.cs.tu-bs.de/courses/ws2223/ag/papers/Ch8/1982-Itai_GridGraphs.pdf
% https://doi.org/10.1137/0211056
\bibitem{Itai1982} A. Itai, C.~H. Papadimitriou, and J.~L. Szwarcfiter. Hamilton paths in grid graphs. \textit{SIAM Journal on Computing}, 11(4):676--686, 1982 \href{https://doi.org/10.1137/0211056}{https://doi.org/10.1137/0211056}.
% https://api.pageplace.de/preview/DT0400.9780080933849_A23543816/preview-9780080933849_A23543816.pdf
% https://api.semanticscholar.org/CorpusID:117976181
\bibitem{Chvatal1985} V. Chvátal. Hamiltonian cycles. In \textit{Handbook of Combinatorics}. Academic Press, 1985 \href{https://api.semanticscholar.org/CorpusID:117976181}{https://api.semanticscholar.org/CorpusID:117976181}.
% https://arxiv.org/pdf/1708.07902
% https://doi.org/10.48550/arXiv.1708.07902
\bibitem{Justesen2019} N. Justesen, P. Bontrager, J. Togelius, and S. Risi. Deep learning for video game playing. \textit{IEEE Transactions on Games}, 12(1):1--20, 2019 \href{https://doi.org/10.48550/arXiv.1708.07902}{https://doi.org/10.48550/arXiv.1708.07902}.
% http://julian.togelius.com/papers/Togelius2009Super.pdf
% https://doi.org/10.1109/CIG.2009.5286481
\bibitem{Togelius2009} J. Togelius, S. Karakovskiy, and J. Koutník. The 2009 two-player {GVGAI} competition. In \textit{IEEE Computational Intelligence and Games (CIG)}, 2009 \href{https://doi.org/10.1109/CIG.2009.5286481}{https://doi.org/10.1109/CIG.2009.5286481}.
% https://github.com/twanvl/snake
% \bibitem{snakeRepo} twanvl. \textit{snake}. GitHub repository, abgerufen am 13.07.2026. \texttt{github.com/twanvl/snake}.
% https://github.com/BrianHaidet/AlphaPhoenix/tree/master/Snake_AI_(2020a)_DHCR_with_strategy
\bibitem{alphaPhoenix} Brian Haidet. \textit{AlphaPhoenix / Snake\_AI\_(2020a)\_DHCR\_with\_strategy}. GitHub repository, abgerufen am 13.07.2026. \href{https://github.com/BrianHaidet/AlphaPhoenix/tree/master/Snake_AI_(2020a)_DHCR_with_strategy}{https://github.com/BrianHaidet/AlphaPhoenix/ tree/master/Snake\_AI\_(2020a)\_DHCR\_with\_strategy}
\end{thebibliography}

\end{document}
